\documentclass{amsart}
% For dealing with references we use the comment environment
\usepackage{verbatim}
\newenvironment{reference}{\comment}{\endcomment}
%\newenvironment{reference}{}{}
\newenvironment{slogan}{\comment}{\endcomment}
\newenvironment{history}{\comment}{\endcomment}

% For commutative diagrams we use Xy-pic
\usepackage[all]{xy}

% We use 2cell for 2-commutative diagrams.
\xyoption{2cell}
\UseAllTwocells

% We use multicol for the list of chapters between chapters
\usepackage{multicol}

% This is generally recommended for better output
\usepackage{lmodern}
\usepackage[T1]{fontenc}

% For cross-file-references

% Package for hypertext links:
\usepackage{hyperref}

% For any local file, say "hello.tex" you want to link to please

% Theorem environments.
%
\theoremstyle{plain}
\newtheorem{theorem}[subsection]{Theorem}
\newtheorem{proposition}[subsection]{Proposition}
\newtheorem{lemma}[subsection]{Lemma}

\theoremstyle{definition}
\newtheorem{definition}[subsection]{Definition}
\newtheorem{example}[subsection]{Example}
\newtheorem{exercise}[subsection]{Exercise}
\newtheorem{situation}[subsection]{Situation}

\theoremstyle{remark}
\newtheorem{remark}[subsection]{Remark}
\newtheorem{remarks}[subsection]{Remarks}

\numberwithin{equation}{subsection}

% Macros
%
\def\lim{\mathop{\mathrm{lim}}\nolimits}
\def\colim{\mathop{\mathrm{colim}}\nolimits}
\def\Spec{\mathop{\mathrm{Spec}}}
\def\Hom{\mathop{\mathrm{Hom}}\nolimits}
\def\Ext{\mathop{\mathrm{Ext}}\nolimits}
\def\SheafHom{\mathop{\mathcal{H}\!\mathit{om}}\nolimits}
\def\SheafExt{\mathop{\mathcal{E}\!\mathit{xt}}\nolimits}
\def\Sch{\mathit{Sch}}
\def\Mor{\mathop{\mathrm{Mor}}\nolimits}
\def\Ob{\mathop{\mathrm{Ob}}\nolimits}
\def\Sh{\mathop{\mathit{Sh}}\nolimits}
\def\NL{\mathop{N\!L}\nolimits}
\def\CH{\mathop{\mathrm{CH}}\nolimits}
\def\proetale{{pro\text{-}\acute{e}tale}}
\def\etale{{\acute{e}tale}}
\def\QCoh{\mathit{QCoh}}
\def\Ker{\mathop{\mathrm{Ker}}}
\def\Im{\mathop{\mathrm{Im}}}
\def\Coker{\mathop{\mathrm{Coker}}}
\def\Coim{\mathop{\mathrm{Coim}}}

% Boxtimes
%
\DeclareMathSymbol{\boxtimes}{\mathbin}{AMSa}{"02}

%
% Macros for moduli stacks/spaces
%
\def\QCohstack{\mathcal{QC}\!\mathit{oh}}
\def\Cohstack{\mathcal{C}\!\mathit{oh}}
\def\Spacesstack{\mathcal{S}\!\mathit{paces}}
\def\Quotfunctor{\mathrm{Quot}}
\def\Hilbfunctor{\mathrm{Hilb}}
\def\Curvesstack{\mathcal{C}\!\mathit{urves}}
\def\Polarizedstack{\mathcal{P}\!\mathit{olarized}}
\def\Complexesstack{\mathcal{C}\!\mathit{omplexes}}
% \Pic is the operator that assigns to X its picard group, usage \Pic(X)
% \Picardstack_{X/B} denotes the Picard stack of X over B
% \Picardfunctor_{X/B} denotes the Picard functor of X over B
\def\Pic{\mathop{\mathrm{Pic}}\nolimits}
\def\Picardstack{\mathcal{P}\!\mathit{ic}}
\def\Picardfunctor{\mathrm{Pic}}
\def\Deformationcategory{\mathcal{D}\!\mathit{ef}}

\setlength{\emergencystretch}{3em}
\begin{document}
\title[Stacks Project, Tag 0AWE]{246 --- Support and Serre S2 property of the canonical module}
\maketitle
\noindent Copyright (C) 2005--2025 Johan de Jong. Stacks Project authors. Source: \href{https://stacks.math.columbia.edu/tag/0AWE}{Tag 0AWE}.

\noindent Editorial difficulty note (not author text). Difficulty H2: The support is identified using normalized localizations, the dimension function and catenarity. The S2 bound is proved by dimension induction, with separate regular-element and depth-zero torsion-removal arguments. The canonical module must satisfy two depth conditions even when the original local ring has depth zero..

\noindent Editorial provenance note (not author text). History: excerpt from revision \texttt{a0444\allowbreak 6e57e\allowbreak c1fbc\allowbreak 252a8\allowbreak 71afc\allowbreak ec775\allowbreak 2fb28\allowbreak 07b14}, dualizing.tex, lines 3358--3412. Reference presentation was adapted; original-author context and core support blocks were retained or added. The mathematical wording is unchanged. Licensed under GNU FDL 1.2 or later; no invariant sections or cover texts. See ../licenses/COPYING. Context blocks reproduce author definitions and supporting results; they are not additional counted problems.

\section{Dualizing complexes over local rings}
\label{section-dualizing-local}

\noindent
In this section $(A, \mathfrak m, \kappa)$ will be a Noetherian local
ring endowed with a dualizing complex $\omega_A^\bullet$ such that
the integer $n$ of Lemma \href{https://stacks.math.columbia.edu/tag/0A7L}{lemma find function (Tag 0A7L)} is zero.
More precisely, we assume that $R\Hom_A(\kappa, \omega_A^\bullet) = \kappa[0]$.
In this case we will say that the dualizing complex is {\it normalized}.
Observe that a normalized dualizing complex is unique up to
isomorphism and that any other dualizing complex for $A$ is isomorphic
to a shift of a normalized one (Lemma \href{https://stacks.math.columbia.edu/tag/0A7F}{lemma dualizing unique (Tag 0A7F)}).



\begin{definition}
\label{algebra-definition-conditions}
Let $R$ be a Noetherian ring.
Let $k \geq 0$ be an integer.
\begin{enumerate}
\item We say $R$ has property {\it $(R_k)$} if for every prime $\mathfrak p$
of height $\leq k$ the local ring $R_{\mathfrak p}$ is regular.
We also say that $R$ is {\it regular in codimension $\leq k$}.
\item We say $R$ has property {\it $(S_k)$} if for every prime $\mathfrak p$
the local ring $R_{\mathfrak p}$ has depth at least
$\min\{k, \dim(R_{\mathfrak p})\}$.
\item Let $M$ be a finite $R$-module. We say $M$ has property $(S_k)$
if for every prime $\mathfrak p$ the module
$M_{\mathfrak p}$ has depth at least
$\min\{k, \dim(\text{Supp}(M_{\mathfrak p}))\}$.
\end{enumerate}
\end{definition}

\begin{lemma}
\label{lemma-depth-dualizing-module}
Let $(A, \mathfrak m, \kappa)$ be a Noetherian local ring with
normalized dualizing complex $\omega_A^\bullet$. Let $d = \dim(A)$
and $\omega_A = H^{-d}(\omega_A^\bullet)$. Then
\begin{enumerate}
\item the support of $\omega_A$ is the union of the irreducible components
of $\Spec(A)$ of dimension $d$,
\item $\omega_A$ satisfies $(S_2)$, see
Algebra, Definition \ref{algebra-definition-conditions}.
\end{enumerate}
\end{lemma}

\begin{proof}
We will use Lemma \href{https://stacks.math.columbia.edu/tag/0A7U}{lemma sitting in degrees (Tag 0A7U)} without further mention.
By Lemma \href{https://stacks.math.columbia.edu/tag/0A7V}{lemma nonvanishing generically local (Tag 0A7V)} the support
of $\omega_A$ contains the irreducible components of dimension $d$.
Let $\mathfrak p \subset A$ be a prime. By Lemma \href{https://stacks.math.columbia.edu/tag/0A7Z}{lemma dimension function (Tag 0A7Z)}
the complex $(\omega_A^\bullet)_{\mathfrak p}[-\dim(A/\mathfrak p)]$
is a normalized dualizing complex for $A_\mathfrak p$. Hence if
$\dim(A/\mathfrak p) + \dim(A_\mathfrak p) < d$, then
$(\omega_A)_\mathfrak p = 0$.
This proves the support of $\omega_A$ is the union of the irreducible
components of dimension $d$, because the complement of this union
is exactly the primes $\mathfrak p$ of $A$ for which
$\dim(A/\mathfrak p) + \dim(A_\mathfrak p) < d$ as $A$ is catenary
(Lemma \href{https://stacks.math.columbia.edu/tag/0A80}{lemma universally catenary (Tag 0A80)}).
On the other hand, if $\dim(A/\mathfrak p) + \dim(A_\mathfrak p) = d$, then
$$
(\omega_A)_\mathfrak p =
H^{-\dim(A_\mathfrak p)}\left(
(\omega_A^\bullet)_{\mathfrak p}[-\dim(A/\mathfrak p)] \right)
$$
Hence in order to prove $\omega_A$ has $(S_2)$ it suffices to show that
the depth of $\omega_A$ is at least $\min(\dim(A), 2)$.
We prove this by induction on $\dim(A)$. The case $\dim(A) = 0$ is
trivial.

\medskip\noindent
Assume $\text{depth}(A) > 0$. Choose a nonzerodivisor $f \in \mathfrak m$
and set $B = A/fA$. Then $\dim(B) = \dim(A) - 1$ and we may apply the
induction hypothesis to $B$. By Lemma \href{https://stacks.math.columbia.edu/tag/0A7T}{lemma divide by nonzerodivisor (Tag 0A7T)}
we see that multiplication by $f$ is injective on $\omega_A$ and we get
$\omega_A/f\omega_A \subset \omega_B$. This proves the depth of $\omega_A$
is at least $1$. If $\dim(A) > 1$, then $\dim(B) > 0$ and $\omega_B$
has depth $ > 0$. Hence $\omega_A$ has depth $> 1$ and we conclude in
this case.

\medskip\noindent
Assume $\dim(A) > 0$ and $\text{depth}(A) = 0$. Let
$I = A[\mathfrak m^\infty]$ and set $B = A/I$. Then $B$ has
depth $\geq 1$ and $\omega_A = \omega_B$ by
Lemma \href{https://stacks.math.columbia.edu/tag/0A7S}{lemma divide by finite length ideal (Tag 0A7S)}.
Since we proved the result for $\omega_B$ above the proof is done.
\end{proof}
\end{document}
